% ============================================================
% SECTION 22 — THE GATEWAY INTERFACE
% ============================================================

\section{The Gateway Interface}

\begin{tcolorbox}[summarybox={\iconFlow[1.5] Abstracting the External World}]
\color{textdark}
In the original monolith, the Orchestrator directly instantiated \code{RazorpayMockService}. This hard-coupling made it impossible to test business logic without triggering payment side-effects, and structurally prohibited the system from easily supporting Stripe, Braintree, or Adyen in the future.

Batch 4 resolved this by introducing the \textbf{Gateway Interface}, a strict abstraction layer enforcing dependency injection across the entire application.
\end{tcolorbox}

\vspace{0.8cm}

% --- THE ABSTRACTION PATTERN ---
\subsection{The Abstraction Pattern}

The application defines a strict Python abstract base class (\code{GatewayInterface}) in \filepath{app/gateways/base.py}. Any payment provider integrated into RecoverAI V2 must implement these exact methods and return standardized \code{ExecutionResult} models.

\begin{center}
\begin{tikzpicture}[
    node distance=1cm and 2cm,
    interface/.style={rectangle, draw=primary, thick, fill=primary!10, minimum width=4cm, minimum height=1cm, rounded corners=4pt, font=\small\bfseries, align=center},
    impl/.style={rectangle, draw=secondary, thick, fill=bglight, minimum width=3cm, minimum height=0.8cm, rounded corners=4pt, font=\scriptsize\ttfamily, drop shadow=black!5, align=center},
    arrow/.style={-{Stealth[scale=1.1]}, thick, draw=textdark, dashed}
]

    \node[interface] (base) {GatewayInterface (ABC)\\ \code{execute\_recovery\_action()}\\ \code{execute\_refund()}};
    
    \node[impl, below left=1.2cm and -1cm of base] (rzp) {RazorpayMockService};
    \node[impl, below=1.2cm of base] (stripe) {StripeService (Planned)};
    \node[impl, below right=1.2cm and -1cm of base] (adyen) {AdyenService (Planned)};

    \draw[arrow] (rzp.north) -- node[left, font=\tiny, sloped] {implements} (base.south west);
    \draw[arrow] (stripe.north) -- (base.south);
    \draw[arrow] (adyen.north) -- node[right, font=\tiny, sloped] {implements} (base.south east);

\end{tikzpicture}
\end{center}

\vspace{0.8cm}

% --- DEPENDENCY INJECTION ---
\subsection{Dependency Injection (\filepath{app/api/dependencies.py})}

To decouple the execution context, the application uses FastAPI's dependency injection system. When \code{ExecutionGuard} or \code{RefundService} needs to talk to the gateway, they do not import Razorpay. They ask the dependency container for "a gateway."

\begin{tcolorbox}[codebox={Injecting the Gateway}]
\begin{lstlisting}[language=Python]
def get_gateway() -> GatewayInterface:
    """
    Returns the configured gateway adapter.
    Currently hardcoded to RazorpayMockService for development.
    """
    return RazorpayMockService()

# Inside ExecutionGuard
class ExecutionGuard:
    def __init__(self, gateway: GatewayInterface):
        self.gateway = gateway
\end{lstlisting}
\end{tcolorbox}

\vspace{0.8cm}

% --- THE MOCK SERVICE ---
\subsection{RazorpayMockService: Simulating Reality}

Because RecoverAI V2 is currently in a pre-production state, the default gateway is \code{RazorpayMockService}. However, it is not a simplistic \code{return True} mock. It is designed to aggressively simulate the unreliability of real banking networks.

\begin{tcolorbox}[featurebox={Simulated Network Behaviors}]
\begin{itemize}
    \item \textbf{Latency:} It sleeps for 0.5 to 2.0 seconds via \code{asyncio.sleep()} or \code{time.sleep()} (depending on the async context) to simulate network latency, testing the application's timeout handling.
    \item \textbf{Failure Injection:} For certain transaction amounts (e.g., \code{9999}), it intentionally raises unhandled exceptions to simulate a 502 Bad Gateway or a socket timeout, verifying the application's \state{UNKNOWN} state recovery.
    \item \textbf{Async Webhooks:} While it returns synchronous success/failure for the API call, it intentionally omits final settlement status, forcing the system to rely on external webhooks (or the reconciliation sweep) to transition to \state{SUCCEEDED}.
\end{itemize}
\end{tcolorbox}

\vspace{0.8cm}

% --- MULTI-TENANCY FOUNDATION ---
\subsection{Foundation for Multi-Tenancy}

The abstraction layer paves the way for \textbf{Multi-Tenancy}. In the future, \code{Transaction} records will include a \code{merchant\_id}. The \code{get\_gateway()} dependency can be upgraded to read the \code{merchant\_id}, look up the merchant's configured provider, and dynamically instantiate \code{StripeService} for Merchant A, and \code{RazorpayService} for Merchant B, without altering a single line of the Orchestrator, Policy Engine, or Execution Guard.
